\subsection{Reweighted Distribution Matching for Dataset Distillation} \label{sec:03-theory-rdm}
\begin{foldable} \label{sec:03-theory-rdm-dm} % DD --> DM --> RDM
    The dataset distillation~(DD) objective seeks $\tilde{\mathcal{D}}$ with $|\tilde{\mathcal{D}}| = K \ll N$ minimizing the loss on the original data:
    \begin{equation}
        \min_{\tilde{\mathcal{D}}}\; \mathbb{E}_{\theta_0 \sim p(\theta_0)} \left[ \mathbb{E}_{z \sim \mathcal{D}} \left[ \ell\!\left(z;\; F(\theta_0;\, \tilde{\mathcal{D}}, \eta, T)\right) \right] \right]
        \label{eq:dataset-distillation}
    \end{equation}
    where $F(\theta_0;\, \tilde{\mathcal{D}}, \eta, T)$ denotes the model obtained by running $T$ steps of gradient descent with learning rate $\eta$ on $\tilde{\mathcal{D}}$ from initialization $\theta_0$.
    This bi-level problem is intractable for discrete text.
    We relax it to distribution matching via the gradient-field equivalence argument~\cite{zhao2023dataset}: if $\tilde{\mathcal{D}}$ induces the same expected gradient field as $\mathcal{D}$ throughout training, the two runs converge to comparable parameters.
    Since $\nabla_\theta \mathbb{E}_{z \sim P}[\ell(z;\theta)] = \mathbb{E}_{z \sim P}[\nabla_\theta \ell(z;\theta)]$, matching distributions in a task-relevant feature space controls the gradient field --- making distribution matching a tractable surrogate for~\eqref{eq:dataset-distillation}.
    This yields the distribution matching~(DM) objective:
    \begin{equation}
        \min_{|\tilde{\mathcal{D}}| = K}\; d\!\left(P_{\tilde{\mathcal{D}}},\; P_{\mathcal{D}}\right)
        \label{eq:distribution-matching}
    \end{equation}
    for a divergence $d$ in a task-relevant feature space.
    Eq.~\eqref{eq:distribution-matching} is a surrogate, not an equivalence; the formal gap bound appears in Appendix~A.3 once $d$ is specified.
\end{foldable}

\begin{foldable} \label{sec:03-theory-rdm-reweight} % Uniform DM --> RDM
    At extreme compression ($\rho = K/N \ll 1$), uniform matching allocates prototypes proportional to density --- but not all regions contribute equally to learning.
    Dense easy regions consume budget that would be better spent on moderate and boundary samples near decision boundaries, which carry the most information for model robustness.
    We instead reweight the source by per-sample importance: define $w_n \geq 0$, $\sum_n w_n = 1$, inducing $P^w = \sum_n w_n \delta_{z_n}$, yielding the reweighted distribution matching~(RDM) objective:
    \begin{equation}
        \min_{|\tilde{\mathcal{D}}| = K}\; d\!\left(P_{\tilde{\mathcal{D}}},\; P^w\right)
        \label{eq:reweighted-distribution-matching}
    \end{equation}
    The weights $w_n$ must reflect downstream learning contribution, not geometry alone, motivating gradient-based estimation.
    Influence self-scores (\S\ref{sec:02-related-influence}) are the natural candidate: they are task-aligned by construction and theoretically grounded.
    We now show that the simplest instantiation---$w_n \propto \mathcal{I}_\text{self}(z_n)$ at a single checkpoint, select top-$K$---fails on both counts: the weights are biased toward noisy samples (Proposition~\ref{prop:hsb}), and independent top-$K$ selection ignores distributional coverage (Remark~\ref{rem:coverage}).
    \begin{remark}[Double relaxation gap] \label{rem:double-gap}
        Moving from~\eqref{eq:dataset-distillation} to~\eqref{eq:distribution-matching} incurs a gradient-field approximation error (bounded formally in Appendix~A.3); moving from~\eqref{eq:distribution-matching} to~\eqref{eq:reweighted-distribution-matching} replaces the uniform source $P_\mathcal{D}$ with the reweighted source $P^w$, incurring an additional cost:
        \begin{equation}
            \mathrm{Gap}_2 \;\leq\; L \cdot W_1(P_\mathcal{D},\, P^w) \;=\; L \cdot W_1\!\left(\tfrac{1}{N}\textstyle\sum_n \delta_{z_n},\; \textstyle\sum_n w_n \delta_{z_n}\right).
            \label{eq:gap2}
        \end{equation}
        Since both measures have the same support, $W_1(P_\mathcal{D}, P^w) \leq \mathrm{diam}(\mathcal{Z}) \cdot \frac{1}{2}\sum_n |w_n - 1/N|$, which is directly computable from the $\kappa_n$ scores and small when $w_n \approx 1/N$ (uniform weights recover DM).
        The reweighting is beneficial precisely when this gap is outweighed by the reduction in downstream loss from correcting the HSB, as formalised in Proposition~\ref{prop:traj} (\S\ref{sec:04-method-gradient}).
    \end{remark}
\end{foldable}


\subsection{The Naive Instantiation Fails: Two Open Problems} \label{sec:03-theory-failures}
\begin{foldable} % Setup
    Define the influence self-score at checkpoint $\hat{\theta}$ as
    \begin{equation}
        \mathcal{I}_\text{self}(z_n; \hat{\theta}) = \nabla_\theta \ell(z_n; \hat{\theta})^\top H_{\hat{\theta}}^{-1} \nabla_\theta \ell(z_n; \hat{\theta}),
        \label{eq:influence-self}
    \end{equation}
    where $H_{\hat{\theta}} = \frac{1}{N}\sum_n \nabla^2_\theta \ell(z_n; \hat{\theta})$ is the empirical Hessian, specializing~\cite{koh2017understanding} by setting the test point equal to the training point.
    High $\mathcal{I}_\text{self}$ identifies atypical examples: samples that are difficult to fit and exert outsized influence on the model's parameters, which in practice are the noisy ones.
    This suggests a naive pipeline: set $w_n \propto \mathcal{I}_\text{self}(z_n; \hat{\theta})$ and return $\arg\max_{|\mathcal{A}|=K} \sum_{i \in \mathcal{A}} w_i$; we identify two orthogonal failures of this strategy.
\end{foldable}

\begin{foldable} % Proposition 1: HSB
    The first failure is a bias in the weights themselves.
    \begin{proposition}[Hard-Sample Bias] \label{prop:hsb}
        Let $\mathcal{D}_\text{clean},\, \mathcal{D}_\text{noisy}$ partition $\mathcal{D}$.
        Assume: \emph{(i)} $\|\nabla_\theta \ell(z; \hat{\theta})\| \leq \epsilon$ for $z \in \mathcal{D}_\text{clean}$;
        \emph{(ii)} $\|\nabla_\theta \ell(z; \hat{\theta})\| \geq \delta \gg \epsilon$ for $z \in \mathcal{D}_\text{noisy}$.
        Let $\mathrm{cond}(H) = \lambda_{\max}(H) / \lambda_{\min}(H)$ be the condition number of $H$.
        Then:
        \begin{equation}
            \Delta_\textup{HSB} := \mathbb{E}_{z \sim \mathcal{D}_\textup{noisy}}[\mathcal{I}_\textup{self}(z;\hat{\theta})]
            - \mathbb{E}_{z \sim \mathcal{D}_\textup{clean}}[\mathcal{I}_\textup{self}(z;\hat{\theta})]
            \;\geq\; \frac{\delta^2 - \mathrm{cond}(H_{\hat{\theta}})\,\epsilon^2}{\lambda_{\max}(H_{\hat{\theta}})}
            \;>\; 0
            \label{eq:hsb}
        \end{equation}
    \end{proposition}
    \begin{proof}[Proof sketch (full proof: Appendix~A.1)]
        Since $\mathcal{I}_\text{self}(z) = \nabla^\top H_{\hat\theta}^{-1} \nabla$, spectral bounds on $H_{\hat\theta}^{-1}$ give
        $\|\nabla\|^2/\lambda_{\max}(H_{\hat\theta}) \leq \mathcal{I}_\text{self}(z) \leq \|\nabla\|^2/\lambda_{\min}(H_{\hat\theta})$.
        Applying the lower bound to noisy samples and the upper bound to clean samples yields \eqref{eq:hsb}.
        Positivity requires $\delta/\epsilon > \sqrt{\mathrm{cond}(H_{\hat\theta})}$, which is mild given $\delta \gg \epsilon$.
    \end{proof}
    Because clean gradients vanish at convergence by definition of a local minimum, the bias is \emph{structural}: it cannot be corrected by post-hoc rescaling of single-checkpoint scores, but only by integrating gradient signals across the training trajectory.
    The binary partition simplifies exposition; in practice the bias is monotone in $\|\nabla\|$, as the spectral bounds above show directly.
    \begin{remark}[Gradient norm as special case] \label{rem:grad-norm}
        Setting $H = I$ in \eqref{eq:influence-self} recovers $\mathcal{I}_\text{self}(z) = \|\nabla\|^2$.
        Under this substitution, \eqref{eq:hsb} reduces to $\Delta_\text{HSB} \geq \delta^2 - \epsilon^2 > 0$, which holds without any spectral assumption on $H$.
        The gradient norm used in TAKE (\S\ref{sec:04-method-gradient}) therefore inherits the same structural bias: it is not merely a tractable proxy but satisfies the HSB lower bound directly.
    \end{remark}
\end{foldable}

\begin{foldable} % Remark 1: independent scoring != coverage
    Even with unbiased weights, a second failure remains: top-$K$ selection is blind to distributional coverage.
    \begin{remark}[Independent scoring $\neq$ distributional coverage] \label{rem:coverage}
        Any rank-and-select strategy decomposes as a sum of independent terms, invariant to pairwise distances between selected samples.
        In contrast, $d(P_{\tilde{\mathcal{D}}}, P^w)$ penalises redundancy and rewards coverage.
        Concretely, if $\phi(z_i) \approx \phi(z_j)$ and both have high $w_i, w_j$, independent scoring selects both; a distributional objective assigns nearly the same cost to selecting just one, freeing a slot for an uncovered region.
    \end{remark}
\end{foldable}

\begin{foldable} % Two open problems
    The two failures compound: biased $w_n$ misdirects a population-level solver toward noisy regions; correct $w_n$ fed into an independent selector still produces redundancy.
    They respectively yield two open problems:
    \begin{enumerate}[label=\textbf{(\alph*)}]
        \item \label{open-problem-weight} \textbf{Biased weighting.} Estimate $w_n$ by integrating gradient signals across trajectory $\{\theta_t\}_{t=0}^{T-1}$, reducing $\Delta_\textup{HSB}$---which we will address in \S\ref{sec:04-method-gradient} (formalised as Proposition~\ref{prop:traj}; Appendix~A.2).
        \item \label{open-problem-select} \textbf{Independent selection.} Solve \eqref{eq:reweighted-distribution-matching} directly for a metric $d$ that jointly optimises coverage of $P^w$---which we will address in \S\ref{sec:04-method-ot}.
    \end{enumerate}
\end{foldable}